\section{Wie mache ich Querverweise?}
\label{sec:querverweise}

Ein Querverweis ist ein Verweis innerhalb eines Dokumentes. z.B.:

\begin{itemize}
    \item auf eine \autoref{tab:simpleTable}
    \item auf eine \autoref{fig:dhbw_logo}
    \item auf ein \autoref{cha:einleitung}
    \item auf ein \autoref{sec:glossar}
\end{itemize}

Das funktioniert entweder mit \lstinline|\ref{fig:dhbw_logo}|, was nur die Nummer zurückgibt, oder mit \lstinline|\autoref{fig:dhbw_logo}|, welches anhand des Präfixes \textit{fig} erkennt, dass es sich um eine Abbildung handelt.

Wichtig: Referenziert werden kann nur, was mit einem Label (\lstinline|\label{fig:dhbw_logo}|) versehen ist. Siehe dazu die bisherigen Beispiele in dieser Arbeit.